\documentclass[11pt]{article}

\usepackage[a4paper,margin=1in]{geometry}
\usepackage{hyperref}
\renewcommand{\UrlFont}{\footnotesize}
\usepackage{enumitem}
\usepackage{titlesec}

\titleformat{\section}{\large\bfseries}{\thesection}{1em}{}
\setlist[itemize]{noitemsep, topsep=4pt}
\setlist[enumerate]{noitemsep, topsep=4pt}

\title{KCPC-CS150 \\ Problem Sheet \\ Data Structures I – Disjoint Set Union (DSU)}
\author{}
\date{}

\begin{document}

\maketitle

\noindent
Developed by the KCPC Internal Committee, this problem sheet accompanies the \textit{Data Structures I} workshop.  
Problems are divided into three categories:

\begin{itemize}
    \item Warm-up problems (basic graph connectivity)
    \item Core practice set (standard DSU and MST applications)
    \item Extended homework (advanced graph modelling and DSU techniques)
\end{itemize}

\hrule
\vspace{1em}

% ----------------------------------------------------
\section*{I. Warm-up Problems}

These problems introduce the concept of graph connectivity and connected components.

\begin{enumerate}

\item \textbf{AtCoder Practice2 A — Disjoint Set Union}

(Classic DSU implementation)

\vspace{0.3em}
{\url{https://atcoder.jp/contests/practice2/tasks/practice2_a}}

\end{enumerate}

\vspace{1em}
\hrule
\vspace{1em}

% ----------------------------------------------------
\section*{II. Core Workshop Practice Set}

These problems should be attempted during the workshop.

They focus on using \textbf{Disjoint Set Union (Union-Find)} and understanding graph connectivity and minimum spanning trees.

\begin{enumerate}

\item \textbf{CSES — Road Construction}

(Dynamic connectivity with DSU)

\vspace{0.3em}
{\url{https://cses.fi/problemset/task/1676/}}

\vspace{0.8em}

\item \textbf{CSES — Road Reparation}

(Minimum Spanning Tree — Kruskal)

\vspace{0.3em}
{\url{https://cses.fi/problemset/task/1675/}}

\vspace{0.8em}

\item \textbf{LeetCode 684 — Redundant Connection}

(Cycle detection with DSU)

\vspace{0.3em}
{\url{https://leetcode.com/problems/redundant-connection/}}

\vspace{0.8em}

\item \textbf{Codeforces 977E — Cyclic Components}

(Detecting special connected components)

\vspace{0.3em}
{\url{https://codeforces.com/problemset/problem/977/E}}

\end{enumerate}

\vspace{1em}
\hrule
\vspace{1em}

% ----------------------------------------------------
\section*{III. Extended Homework}

\textbf{Important.} The following problems are significantly more difficult.  
They should only be attempted after completing all problems in the warm-up and core sections.

These exercises introduce more advanced graph modelling and DSU-based techniques.

\medskip

\textbf{Matthew's Note.}  
If you have any questions about the problems or underlying concepts, feel free to contact me or the Director of Competitive Programmes (Marcell).

\vspace{0.8em}

\begin{enumerate}

\item \textbf{CSES — Building Roads}

(Connected components and graph construction)

\vspace{0.3em}
{\small\url{https://cses.fi/problemset/task/1666/}}

\vspace{0.8em}

\item \textbf{Codeforces 1151B — Dima and a Bad XOR}

(Graph reasoning + constructive logic)

\vspace{0.3em}
{\small\url{https://codeforces.com/problemset/problem/1151/B}}

\vspace{0.8em}

\item \textbf{Codeforces 1514D — Cut and Stick}

(Advanced data structures + query processing)

\vspace{0.3em}
{\small\url{https://codeforces.com/problemset/problem/1514/D}}

\vspace{0.8em}

\item \textbf{Codeforces 343D — Water Tree}

(Tree operations + advanced data structures)

\vspace{0.3em}
{\small\url{https://codeforces.com/problemset/problem/343/D}}

\vspace{0.8em}

\item \textbf{Codeforces 500B — New Year Permutation}

(Graph connectivity + permutation reconstruction)

\vspace{0.3em}
{\small\url{https://codeforces.com/problemset/problem/500/B}}

\vspace{0.8em}

\item \textbf{HackerRank — Merging Communities}

(Classic DSU with component size queries)

\vspace{0.3em}
{\small\url{https://www.hackerrank.com/challenges/merging-communities}}

\end{enumerate}

\vspace{1em}

\hrule
\vspace{1em}

\noindent
\textbf{Guidance:}

\begin{itemize}
\item Learn the Disjoint Set Union structure (Union-Find).
\item Use \textbf{path compression} and \textbf{union by rank / size}.
\item Recognise when connectivity queries can be solved with DSU instead of graph traversal.
\item For MST problems, apply \textbf{Kruskal's algorithm}.
\end{itemize}

\end{document}